\chapter{Discussion}\label{chap:discussion}

% 1. Summary
In this thesis we have presented the first model to unify the rich multiscale dynamics of visual search. Previous work has focused their study to models of either large-scale saccades or small-scale fixations. Further, motivations for such studies have been divided in either aiming to present a biologically accurate understanding of visual search, or a full mechanistic generative model of eye-tracking data. Taking aim to unify these motivations and spatiotemporal scales of study, we constructed a model from the neural dynamics developed under the theory of fractional neural sampling (FNS) \cite{qi_fractional_2022, chen_spatiotemporal_2022}. Under this framework, gaze dynamics are governed by a second-order stochastic process (\cref{eq:fns_system}) driven by Lévy noise exploring a dynamical saliency landscape adapts through a dynamical-informative kernel (\cref{sec:dynamical-information_adaptation}). We understand the model dynamics through three key spatiotemporal signatures: global $\alpha$--Lévy superdiffusion, rhythmic transitions between potential wells driven by an adaptation from a kernel function $k(\cdot)$, and local high-frequency harmonic oscillations driven by a memory--momentum parameter $\beta$. Each of these signatures maps onto a well-known behaviour during gaze dynamics. We compared each of these dynamics with the high-resolution gaze data of \textcite{naher_primate_saccades} and with the wider literature, and found that the model captures the signature spatiotemporal dynamics of saccades and fixations within a single framework (\cref{chap:empirical_validations}). In particular, we discovered robust monoscaling Lévy superdiffusion on all spatial scales with $\alpha = 1.12$, theta-rhythmic saccades with $t_{\theta} \approx \qty{200}{\milli \second}$ and a high degree of rhythmicity $\text{RCRI}=0.72$, and local high-frequency harmonic oscillations with mean frequency $f=\qty{56.5}{\hertz}$. Finally, after validating the biologically-motivated model of gaze dynamics, we formulated functional explanations for each of their existence. Using information-theoretic metrics on the updating belief distribution present in the model, we show that each pattern plays a distinct functional role in the efficiency of informative search (\cref{chap:functional_advantages}). 

% 2. The novel physics
The methods and dynamics here studied are of great importance to physicists as a description of a non-equilibrium complex dynamical system. First, visual dynamics are shown to be superdiffusive, and moreover display robust monoscaling, governed by a single stability index $\alpha$. Such anomalous diffusion at every scale is of interest when studying how basic computational dynamics may be constructed from complex dynamical systems at criticality. We characterised the short-time, noise-driven scaling of such biological data through fractional absolute moments (\cref{sec:search_superdiffusion}), and compared these to the monoscaling that an implementation of Lévy dynamics can provide (\cref{sec:levy_superdiffusion}). Such rich multiscale self-similarity has a physical explanation as the result of a complex dynamical network at criticality with diverging correlation length, resulting in nonlocal spatial access to an informative landscape. Second, the model presents and explains the unique pattern of nested oscillations: theta-scale rhythmic transitions between regions, which modulate gamma-scale sampling dynamics, arising from neural adaptation, and the confinement of a sampling process. Nested timescales are a recurrent feature of coherent neural dynamics, especially during communication of information between areas of the brain \cite{tan_theta_2024, spyropoulos_theta_2018, fries_mechanism_2005, bastos_communication_2015, fries_rhythms_2015, fries_rhythmic_2023}. The present model carries this framework over to gaze dynamics, as they emerge from the same neural dynamics which control the process of active vision. Third, the dynamical-informative kernel introduced in this work describes how visual dynamics are influenced by what is perceived, so the searcher depletes information as it samples an informative distribution $\pi$ in search for a target. Our formulation naturally encodes this information through a belief landscape which evolves from a prior saliency $\pi_0$ to a posterior belief after exploration $\pi_t$. The cumulative information gathered is then well defined through the Kullback--Leibler (KL--) divergence. Such a well-motivated physics-inspired development of interpretable memory-effects, heavy-tailed noise, and nonlocality is important in the development and progress of physics as an interdisciplinary field.

% 3. Importance for neuroscience
This work advances the study of the brain in neuroscience, as our model turns gaze into a source of interpretable parameters modulating neural attention, a very important topic of study when understanding decision-making in the brain. The active parameter space here defined through the stability index $\alpha$, memory--momentum strength $\beta$, noise scale $\gamma$, and informative timescale $t_{\theta}$; each parameter has a well-understood and motivated have a defined effect on the trajectory, so fitting them to data estimates the dynamics of the underlying attentional control. Several experiments would sharpen this. High-frequency eye tracking without a bite-bar, combined with head acceleration measurements, would resolve gamma-scale fixational dynamics in naturalistic, head-free viewing and separate eye from head contributions. High-quality smooth pursuit experiments, in which the target moves, would probe how the landscape is tracked over time. Dynamical search tasks with controlled likelihood landscapes would constrain the kernel and the theta-rhythmic transitions. Because gaze aggregates the activity of many regions, the fitted parameters may summarise collective properties of the circuits that control and process visual input. A natural candidate is the interaction between the superior colliculus, which integrates sensory, cortical, and goal-related signals into a map that drives saccades \cite{[CITE]}, and the visual cortex, which supplies the evidence on which that map is built \cite{[CITE]}. Within the active vision loop, the model's evolving potential could be compared with such a priority map, and the kernel with processes such as inhibition of return \cite{ludwig_mechanism_2009, redden_inhibition_2021}, giving a way to study attention in this system quantitatively. Because alterations of eye movements accompany many neurological and psychiatric conditions \cite{[CITE]}, mechanistic parameters of this kind also offer a route to more accurate biomarkers, and potentially to treatments that target a specific part of the loop.

This work advances the study of the brain in neuroscience, as our model turns gaze into a source of interpretable parameters modulating neural attention, a very important topic of study when understanding decision-making in the brain. The active parameter space here defined through the stability index $\alpha$, memory--momentum strength $\beta$, noise scale $\gamma$, and informative timescale $t_{\theta}$; each parameter has a well-understood and motivated effect on the dynamics and function of visual sampling, so fitting them to data estimates the dynamics of the underlying attentional control of the aggregated network of the brain. 

Several experiments may be developed to apply this framework to understand brain health in particular. An experimental protocol tracking high-frequency eye movement alongside coarse-grained saccade positions during a standardised search task without head restraint, combined with head acceleration measurements, would resolve gamma-scale fixational dynamics in naturalistic, head-free viewing and separate eye from head contributions. During such an experiment, coherence of the high-frequency sampling during fixations can be examined rigorously, as well as testing the saliency-frequency coupling and mapping the paramter-space dynamics as described in \cref{sec:spatiotemporal_predictions}. Such an experiment is simple to design using new low-cost and non-invasive techniques such as described in \cite{ryle_simultaneous_2015} and the data extracted is tractable to analyse across the scales of interest here, to test the predictions of FNS dynamics here presented.

High-quality smooth pursuit experiments, tracking a dynamical target, would probe dynamical attention. The present model is easily extended to a dynamical potential, where information adaptation is modelled to forget recent information as it becomes less reliable in a dynamical environment. As such, a potential may be represented as:
\begin{equation}
    V_t(\textbf{x}) = V_{\text{dynamic}}(\textbf{x};t) + K(\textbf{x};t), \qquad \partial_t K(\textbf{x};t) = \frac{1}{\tau}k(\textbf{x}-\textbf{x}_t) - \frac{1}{\tau'} K(\textbf{x};t)
\end{equation}

where sych dynamics will constrain the noise produced, for instance a slow-moving pendulum is predicted to have high rate of large micro-saccades at the apex, whereas a faster pendulum will slow the rate of kernel deposition, making microsaccades as likely as during the smooth pursuit of the pendulum. Because gaze aggregates the activity of many regions, the fitted parameters during any given task will provide an aggregate understanding of the neural circuits that control and process visual input, which are a significant part of the brain. A natural candidate for locating these neural interactions is between the superior colliculus, which integrates sensory, cortical, and goal-related signals into a map that drives saccades \cite{[CITE]}, and the visual cortex, which supplies the evidence on which that map is built \cite{[CITE]}. Because dynamics of eye movements accompany many neurological and psychiatric conditions \cite{[CITE]}, mechanistic parameter fitting of this kind also offer a route to highly accurate, simple, and non-invasive biomarkers.

% 4. Importance and application to machine learning
The same ideas apply to machine learning, in particular to sampling from multimodal and heavy-tailed distributions. Standard Markov chain Monte Carlo methods based on local moves struggle to cross between well-separated modes. The model's superdiffusive jumps address this directly, and its adaptive kernel resembles well-tempered metadynamics \cite{tiwary_time-independent_2015}, in which deposited bias discourages revisiting explored regions. The link to free-energy minimisation and active inference \cite{NIPS2017_194cf6c2, [CITE]} is also close, since both treat search as the reduction of uncertainty about a landscape. Our results emphasise that coupling oscillations across scales is itself important for an abstract sampler: rapid local oscillations explore each mode, while slower rhythmic transitions decide when to move on. Superdiffusive samplers have recently attracted considerable interest \cite{[CITE]}. The same principle has influenced more general AI frameworks, notably in work from our group on fractional neural attention (FNA) \cite{qu_frac}. FNA models token interactions through L\'evy diffusion governed by the fractional Laplacian, so that short- and long-range dependencies arise from a single operator rather than hand-designed local--global patterns. In the reported experiments it matches or exceeds Transformer baselines on text classification, image classification, and machine translation, in some cases with a single layer, and its attention networks show larger spectral gaps and shorter path lengths, consistent with faster information mixing. That the same mechanism describes both gaze and attention in artificial networks supports the view that multiscale, nonlocal sampling is a general computational strategy.

% 5. Broad conclusions
In conclusion, this thesis has shown that the rich and apparently disparate dynamics of human gaze can be described by a single physical model with a small set of interpretable components which derive from neural dynamics: nonlocal Lévy jumps, nested oscillations, and an adaptive, information-gathering search landscape. Gaze is both easily measured and serves as a strong proxy for neural decision-making and attention, so gaining an understanding of gaze dynamics gives access to basic computation in the brain. The model offers a quantitative framework making use of concepts from the statistical physics of diffusive processes, makes testable predictions for new experiments to test the theory, and provides a process through which we can understand neural attention at a large scale, and interpret it to develop and understand brain health and basic function of abstract dynamical computational systems. 